% ─────────────────────────────────────────────────────────────────────────────
% Deliverable 3 — Test Plan
% MATH5320 Portfolio Risk Management System
% Columbia University, MATH GR 5320 Financial Risk Management, Spring 2026
% ─────────────────────────────────────────────────────────────────────────────
\documentclass[11pt]{article}
\usepackage[margin=1in]{geometry}
\usepackage[T1]{fontenc}
\usepackage[utf8]{inputenc}
\usepackage{lmodern}
\usepackage{microtype}
\usepackage{xcolor}
\usepackage{amsmath,amssymb}
\usepackage{booktabs,longtable,array,ragged2e}
\usepackage{graphicx}
\usepackage{float}
\usepackage{hyperref}
\usepackage{enumitem}
\usepackage{parskip}
\usepackage{fancyvrb}

\hypersetup{
  colorlinks=true,
  linkcolor=blue,
  urlcolor=blue,
  citecolor=blue,
  pdftitle={Test Plan — MATH5320},
  pdfauthor={Nigel Li, Michael Adegbite, Stella}
}

\setlist[itemize]{topsep=3pt,itemsep=2pt,parsep=0pt,leftmargin=1.6em}
\setlist[enumerate]{topsep=3pt,itemsep=2pt,parsep=0pt,leftmargin=1.6em}
\newcolumntype{L}[1]{>{\RaggedRight\arraybackslash}p{#1}}
\newcommand{\code}[1]{\texttt{#1}}
\DefineVerbatimEnvironment{shellcode}{Verbatim}{fontsize=\small,xleftmargin=1em}

% ─────────────────────────────────────────────────────────────────────────────
\begin{document}

% ── Cover page ────────────────────────────────────────────────────────────────
\begin{titlepage}
\centering
\vspace*{2cm}
{\LARGE\bfseries Test Plan\par}
\vspace{0.4cm}
{\large MATH GR 5320 Portfolio Risk Management System\par}
\vspace{0.4cm}
{\large Columbia University — Financial Risk Management, Spring 2026\par}
\vspace{1.4cm}

\begin{tabular}{ll}
\toprule
\textbf{Field} & \textbf{Value} \\
\midrule
Deliverable & 3 of 5 (20 points) \\
Authors & Nigel Li, Michael Adegbite, Stella \\
Reference Commit & \texttt{5841589e3f3d2dbd3c1e38b08642eccce201a6a2} \\
Submission Date & May 12, 2026 \\
Python Version & 3.12.2 \\
Test Suite & 610 tests, 0 failures (see Deliverable~5) \\
Statement Coverage & 96\% \\
\bottomrule
\end{tabular}

\vspace{1.4cm}
\begin{minipage}{0.78\linewidth}
\small\itshape
This test plan formalizes the validation program for the MATH5320 risk
engine.  It treats testing as part of model development and model
governance, aligned with Lecture~5's emphasis on pre-deployment
model risk management.
\end{minipage}
\end{titlepage}

\tableofcontents
\clearpage

% ─────────────────────────────────────────────────────────────────────────────
\section{Executive Summary}

This test plan addresses four questions that every model validation
program must answer:

\begin{enumerate}
  \item What is being tested?
  \item Why is it being tested?
  \item Against what benchmark is it being tested?
  \item What constitutes acceptable behavior?
\end{enumerate}

The \code{MATH5320} repository contains a broad test suite covering
formula correctness, portfolio valuation, VaR/ES calculations,
backtesting, market-data loading, UI behavior, and course-formula
extensions.  The plan below formalizes those tests into a
model-risk-oriented validation program.  Testing is treated as a
first-class component of model development and model governance, not
merely as proof that the application executes without crashing.

The plan also includes a dedicated benchmark-comparison section
(Section~11) that evaluates our implementations against external
analytical references — a requirement highlighted in the Lecture~5
model validation framework.

% ─────────────────────────────────────────────────────────────────────────────
\section{Test Objectives}

The test plan validates the following properties:

\begin{enumerate}
  \item Correctness of mathematical formulas;
  \item Correctness of portfolio valuation and option repricing;
  \item Correctness of VaR/ES calculations under historical, parametric, and Monte Carlo methods;
  \item Correctness of backtesting and exception logic;
  \item Correctness of credit and regulatory formula-sheet extensions;
  \item Correct handling of invalid inputs, edge cases, and numerical failure modes;
  \item Correct service-layer integration and UI behaviour;
  \item Reproducibility through deterministic fixtures, seeded Monte Carlo, and coverage reporting.
\end{enumerate}

These objectives align with Lecture~5's emphasis on testing as part of
model development and validation, rather than mere application execution.

% ─────────────────────────────────────────────────────────────────────────────
\section{Scope of Testing}

\textbf{In scope:}
pure pricing formulas; portfolio valuation and aggregation; return and
covariance estimation; historical VaR/ES; parametric VaR/ES; Monte Carlo
VaR/ES; VaR backtesting and exception diagnostics; course-formula extensions
in \code{src/risk/lognormal.py}, \code{src/credit/}, and
\code{src/risk/regulatory.py}; market-data loading and validation; and
Streamlit UI rendering.

\textbf{Out of scope:}
production deployment hardening; enterprise access control; performance
benchmarking at scale; full external market-data certification; and
production volatility-surface validation.

% ─────────────────────────────────────────────────────────────────────────────
\section{Test Environment}

\begin{itemize}
  \item Repository root: \code{MATH5320}
  \item Python: \code{3.12.2}; OS: macOS / Darwin arm64
  \item Key packages: \code{streamlit 1.37.1}, \code{numpy 1.26.4},
    \code{pandas 3.0.2}, \code{scipy 1.17.1}, \code{plotly 5.24.1},
    \code{yfinance 1.2.0}, \code{pytest 7.4.4}, \code{pytest-cov 7.1.0}
\end{itemize}

Environment snapshot commands:
\begin{shellcode}
git rev-parse HEAD
python --version
pip freeze > test_artifacts/requirements_freeze.txt
\end{shellcode}

Network requirements: the no-network unit suite runs without external
downloads.  Live-data integration scripts require network access and are
explicitly separated.  All Monte Carlo regression tests use fixed seeds;
if a seed is randomized for exploratory work, it is logged.

% ─────────────────────────────────────────────────────────────────────────────
\section{Test Data and Fixtures}

\subsection{Synthetic Fixtures}

Used for deterministic unit tests: synthetic price histories, toy
two-stock portfolios, simple option positions, and deterministic
exposures and covariance matrices.

\subsection{Course-Derived Fixtures}

Used for validation against homework and formula-sheet results:
exact GBM/lognormal values, hazard-rate survival and spread values,
Merton Q/P default probabilities and valuations, CDS and CVA examples,
and regulatory RWA and capital-ratio examples.

\subsection{Bloomberg Course Data}

Observed local files: \code{data/AAPL-bloomberg.csv} and
\code{data/CAT-bloomberg.csv}.  These support the AAPL/CAT course
portfolio notebooks and provide course-accepted regression anchors.

\subsection{Live Market Data}

Used by integration scripts: Yahoo Finance downloads for equities and
rate proxies via the cached download layer in \code{src/data/market\_data.py}.

\subsection{Data Proxy Policy}

Consistent with Lecture~5's requirement to document data quality, proxies,
and cleaning assumptions:

\begin{itemize}
  \item CSV-based Bloomberg fixtures are treated as the primary course
    acceptance data.
  \item Yahoo Finance is treated as a convenience data source, not a
    gold-standard benchmark.
  \item Any cleaning, alignment, or row-dropping is performed explicitly
    and logged.
\end{itemize}

% ─────────────────────────────────────────────────────────────────────────────
\section{Test Categories}

\begin{center}
\begin{tabular}{L{3.5cm}L{5cm}L{5cm}}
\toprule
\textbf{Category} & \textbf{Purpose} & \textbf{Example} \\
\midrule
Unit tests              & Validate pure functions & Black-Scholes price, hazard survival \\
Analytical goldens      & Compare against closed-form values & Exact GBM VaR/ES \\
Homework fixtures       & Course-derived regression values & AAPL/CAT VaR, Merton Q/P \\
External benchmarks     & Independent comparison source & Option-calculator Black-Scholes, Basel traffic light \\
Edge cases              & Validate boundary behaviour & Zero hazard, zero volatility \\
Failure-mode tests      & Ensure controlled errors & Invalid confidence, missing data \\
Behavioural tests       & Financial monotonicity and logic & Option price increases with vol \\
Convergence tests       & Numerical stability & MC VaR converging as simulation count increases \\
Backtesting tests       & VaR forecast logic & Exceptions, Kupiec LR, Christoffersen LR \\
Data validation tests   & Input quality checking & Missing or stale prices \\
Integration tests       & Full workflow & Portfolio input to risk output \\
UI tests                & Streamlit panel behavior & Portfolio editor, settings, results \\
Coverage tests          & Source execution breadth & Coverage report and missing-line review \\
\bottomrule
\end{tabular}
\end{center}

% ─────────────────────────────────────────────────────────────────────────────
\section{Module-Level Test Matrix}

\begin{center}
\begin{tabular}{L{5cm}L{8.5cm}}
\toprule
\textbf{Module} & \textbf{Required tests} \\
\midrule
\code{pricing/black\_scholes.py}          & Price, delta, put-call parity, invalid inputs, monotonicity \\
\code{portfolio/positions.py}             & Stock value, option value, delta exposure, long/short signs \\
\code{portfolio/portfolio.py}             & Aggregate value, exposure vector, empty portfolio rejection \\
\code{risk/returns.py}                    & Log returns, overlapping returns, horizon summation \\
\code{risk/estimators.py}                 & Rolling mean/cov, EWMA, covariance symmetry \\
\code{risk/historical.py}                 & Historical VaR/ES, log shock, missing data \\
\code{risk/parametric.py}                 & Normal VaR/ES, covariance aggregation, ES confidence separation \\
\code{risk/monte\_carlo.py}               & Seeded reproducibility, MC VaR/ES, covariance validation \\
\code{risk/backtest.py}                   & Exceptions, Kupiec, Christoffersen, traffic light, severity \\
\code{risk/lognormal.py}                  & Exact long/short GBM VaR/ES \\
\code{credit/hazard.py}                   & Constant and piecewise hazard, survival, density, risky ZCB \\
\code{credit/merton.py}                   & Q/P PD, equity/debt, target-survival inversion \\
\code{credit/cds.py}                      & Approximation and full par spread \\
\code{credit/cva.py}                      & EPE, CVA, discounted CVA \\
\code{credit/mitigation.py}               & Netting, collateral, CSA logic \\
\code{risk/regulatory.py}                 & RWA, capital ratio, DFAST pathing \\
\code{services/risk\_engine\_service.py}  & Orchestration and result-object consistency \\
\code{ui/*.py}                            & Streamlit input handling and result rendering \\
\bottomrule
\end{tabular}
\end{center}

% ─────────────────────────────────────────────────────────────────────────────
\section{Analytical Golden Tests}

\subsection{Black-Scholes}

Required test family:

\begin{itemize}
  \item \textbf{BS\_01} call price against known analytical value
  \item \textbf{BS\_02} put price against known analytical value
  \item \textbf{BS\_03} put-call parity: $C - P = S e^{-qT} - K e^{-rT}$
  \item \textbf{BS\_04} call delta $\in [0,1]$
  \item \textbf{BS\_05} put delta $\in [-1,0]$
  \item \textbf{BS\_06} option price increases with volatility (vega $> 0$)
  \item \textbf{BS\_07} invalid maturity or volatility raises controlled exception
\end{itemize}

Acceptance criterion: numerical values agree within analytic tolerance;
parity holds within floating-point tolerance; invalid domains fail loudly.

\subsection{Exact Lognormal VaR/ES}

Required test family:

\begin{itemize}
  \item \textbf{LN\_01} long VaR: $V_0[1 - e^{m_h + s_h z_{1-\alpha}}]$
  \item \textbf{LN\_02} long ES: closed-form formula with $\mathcal{N}(z_{1-\alpha} - s_h)/(1-\alpha)$
  \item \textbf{LN\_03} short VaR: $V_0[e^{\mu h + z_\alpha \sigma\sqrt{h}} - 1]$
  \item \textbf{LN\_04} short ES: closed-form short formula
  \item \textbf{LN\_05} VaR scales linearly with notional
  \item \textbf{LN\_06} short VaR exceeds long VaR for identical base inputs
  \item \textbf{LN\_07} zero horizon gives zero-risk limit or controlled rejection
\end{itemize}

\subsection{Normal Parametric VaR/ES}

Required test family:

\begin{itemize}
  \item \textbf{NORM\_01} VaR formula: $-\mu_h^\top x + \sqrt{x^\top \Sigma_h x} \cdot \Phi^{-1}(\alpha)$
  \item \textbf{NORM\_02} ES formula: analogous closed-form integral
  \item \textbf{NORM\_03} ES $\geq$ VaR when using the same confidence level
  \item \textbf{NORM\_04} covariance aggregation for multi-asset portfolio
  \item \textbf{NORM\_05} perfectly offsetting exposures reduce risk
  \item \textbf{NORM\_06} invalid covariance is rejected or handled explicitly
\end{itemize}

% ─────────────────────────────────────────────────────────────────────────────
\section{Homework-Derived Regression Fixtures}

The repository contains substantial homework-derived cases that function
as formal regression fixtures.  These provide direct numerical anchors —
exactly the kind of evidence Lecture~5 encourages.

\begin{center}
\begin{tabular}{L{4.5cm}L{3cm}L{6cm}}
\toprule
\textbf{Case ID} & \textbf{Area} & \textbf{Expected validation} \\
\midrule
\code{HW4\_SINGLE\_STOCK}    & GBM VaR                   & 5-day 99\% VaR near homework value \\
\code{HW4\_TWO\_STOCK}       & Parametric covariance VaR & Correct mean/variance/correlation aggregation \\
\code{HW6\_EWMA}             & Rolling/EWMA estimation   & Window and EWMA parameter behavior \\
\code{HW6\_HAZARD\_CONST}    & Constant hazard           & Survival/default probability regression \\
\code{HW6\_HAZARD\_PIECEWISE}& Piecewise hazard          & $\lambda(t)$, $\Lambda(t)$, $s(t)$, $p(t)$, spread table \\
\code{HW7\_MERTON\_QP}       & Merton                    & Q vs P PD comparison \\
\code{HW7\_MERTON\_TIMING}   & Merton timing             & Zero default probability before maturity interval \\
\code{HW8\_CDS}              & CDS                       & Constant-hazard approx spread and full par spread \\
\code{HW8\_CVA}              & CVA                       & Exposure and default-probability aggregation \\
\code{HW9\_MERTON\_INVERSION}& Merton inversion          & Target-survival inversion \\
\code{HW9\_SHORT}            & Short-risk formulas       & Short VaR/ES sign and magnitude behavior \\
\code{HW10\_RWA}             & Regulatory                & RWA and capital-ratio arithmetic \\
\code{HW10\_DFAST}           & Regulatory                & 9-quarter stress-path structure \\
\bottomrule
\end{tabular}
\end{center}

% ─────────────────────────────────────────────────────────────────────────────
\section{External Benchmark Tests}

External benchmarks are not substitutes for homework fixtures.
They provide a second independent source of validation confidence.

\subsection{Option Calculator Benchmark Cases}

\begin{center}
\begin{tabular}{lrrrrrrr}
\toprule
\textbf{Case} & $S$ & $K$ & $T$ & $r$ & $q$ & $\sigma$ & \textbf{Type} \\
\midrule
ATM call    & 100 & 100 & 1.0 & 0.05 & 0.00 & 0.20 & call \\
ATM put     & 100 & 100 & 1.0 & 0.05 & 0.00 & 0.20 & put \\
Dividend call & 100 & 105 & 2.0 & 0.03 & 0.02 & 0.25 & call \\
ITM put     &  90 & 100 & 0.5 & 0.04 & 0.01 & 0.30 & put \\
Near-expiry & 100 & 100 & $1/252$ & 0.05 & 0.00 & 0.20 & call \\
\bottomrule
\end{tabular}
\end{center}

Acceptance criterion: price and delta agree within standard numerical
rounding tolerance of an external Black-Scholes calculator.

\subsection{Basel Traffic-Light Benchmark}

If the Basel-zone helper is present, required assertions are:

\begin{itemize}
  \item \code{basel\_zone(0, 250, 0.99) == green}
  \item \code{basel\_zone(4, 250, 0.99) == green}
  \item \code{basel\_zone(5, 250, 0.99) == amber}
  \item \code{basel\_zone(9, 250, 0.99) == amber}
  \item \code{basel\_zone(10, 250, 0.99) == red}
\end{itemize}

\subsection{DFAST Structural Benchmark}

We do not claim Federal Reserve replication.  The structural test checks:
9-quarter stress path; baseline / adverse / severely adverse scenario
naming; Tier~1 capital path; RWA path; and minimum capital ratio threshold.

% ─────────────────────────────────────────────────────────────────────────────
\section{MRM Benchmark Model Comparison}
\label{sec:mrm-benchmark}

A model validation program following Lecture~5's framework requires
comparison of the model under review against an independent reference or
benchmark approach.  We perform three benchmark comparisons.

\subsection{Parametric vs.\ Historical Benchmark}

The parametric delta-normal model serves as the primary benchmark for
the historical simulation model, and vice versa.  For any given
portfolio and data window, we verify:

\begin{itemize}
  \item Both models produce positive VaR and ES.
  \item ES $\geq$ VaR at the same confidence level for both methods.
  \item For a simple single-stock equity-only portfolio (no options),
    parametric and historical VaR agree within a pre-specified tolerance
    of approximately 10\%.  Larger differences are expected and explicitly
    documented when the historical return distribution departs significantly
    from normality.
  \item Adding options consistently increases both historical and MC VaR
    relative to the corresponding linear (parametric) estimate when
    significant nonlinearity is present.
\end{itemize}

This comparison is implemented in \code{tests/test\_backend.py} and
\code{tests/test\_course\_validation.py}, and the results are recorded in
\code{submission/test\_artifacts/official\_benchmark\_results.csv}.

\subsection{Monte Carlo vs.\ Exact GBM Benchmark}

For a single-asset position with no options, Monte Carlo VaR should
converge to the exact GBM/lognormal closed-form VaR as the simulation
count increases.  We verify:

\begin{itemize}
  \item At $n_{\mathrm{sims}} = 100\,000$ with a fixed seed, Monte Carlo
    99\% VaR agrees with \code{var\_long\_lognormal} within 2\% for a
    representative single-stock case.
  \item Convergence is monotone in expectation as $n_{\mathrm{sims}}$ grows.
\end{itemize}

This benchmark is implemented in \code{tests/test\_lognormal.py} and
\code{tests/test\_coverage\_gaps.py}.

\subsection{CDS Approximation vs.\ Full Par Spread Benchmark}

The constant-hazard CDS par-spread approximation
$s_{\mathrm{CDS}} \approx (1-R)\lambda$ is benchmarked against the
full discrete-payment formula from \code{cds\_par\_spread}.  We verify:

\begin{itemize}
  \item For $\lambda = 3\%$ and $R = 40\%$, the approximation yields
    approximately $180$ basis points, consistent with the §14 course landmark.
  \item The full formula produces a result within 5\% of the approximation
    for short tenors and flat hazard curves.
  \item The full formula diverges predictably from the approximation
    as the tenor increases or the hazard rate increases, confirming
    that the approximation error is understood and bounded.
\end{itemize}

This benchmark is implemented in \code{tests/test\_credit.py} and
\code{tests/test\_course\_validation.py}.

% ─────────────────────────────────────────────────────────────────────────────
\section{Edge-Case and Failure-Mode Tests}

The following edge and failure cases are explicitly required and all
are implemented in the test suite:

\begin{itemize}
  \item \textbf{EDGE\_01} Empty portfolio raises a controlled exception
  \item \textbf{EDGE\_02} Missing ticker history raises a controlled exception
  \item \textbf{EDGE\_03} Insufficient lookback raises with a descriptive message
  \item \textbf{EDGE\_04} NaN prices are handled or rejected explicitly
  \item \textbf{EDGE\_05} Duplicate dates are handled or rejected explicitly
  \item \textbf{EDGE\_06} Negative or zero prices are rejected
  \item \textbf{EDGE\_07} Invalid confidence level is rejected
  \item \textbf{EDGE\_08} VaR confidence can be set independently from ES confidence
  \item \textbf{EDGE\_09} Non-positive volatility is rejected
  \item \textbf{EDGE\_10} Zero maturity is handled or rejected
  \item \textbf{EDGE\_11} Non-PSD covariance is handled or rejected
  \item \textbf{EDGE\_12} Monte Carlo seed reproducibility is confirmed
  \item \textbf{EDGE\_13} $n_{\mathrm{sims}} \leq 0$ is rejected
  \item \textbf{EDGE\_14} Zero hazard gives survival $= 1$ and PD $= 0$
  \item \textbf{EDGE\_15} Recovery $R = 1$ gives zero CDS/CVA loss
  \item \textbf{EDGE\_16} Merton survival decreases as debt face value increases
  \item \textbf{EDGE\_17} Capital-ratio division by zero is rejected
  \item \textbf{EDGE\_18} Netted exposure $\leq$ gross exposure
  \item \textbf{EDGE\_19} Collateralized exposure $\leq$ uncollateralized exposure
  \item \textbf{EDGE\_20} ES $\geq$ VaR when evaluated at the same confidence
\end{itemize}

% ─────────────────────────────────────────────────────────────────────────────
\section{Integration and UI Tests}

\subsection{Integration Tests}

Required integration paths:

\begin{itemize}
  \item Portfolio creation to \code{RiskEngineService.run\_all()};
  \item End-to-end VaR/ES computation under live or cached market data;
  \item End-to-end backtesting with Kupiec and Christoffersen diagnostics;
  \item Formula-sheet integration with live market data.
\end{itemize}

Relevant files: \code{tests/integration\_test.py} and
\code{tests/integration\_test\_formula\_sheet.py}.

\subsection{UI Tests}

Required UI test areas: portfolio editor; market data panel; risk
settings; results panel; credit panel; CDS/CVA panel; capital panel;
and chart helpers.

Relevant files: \code{tests/test\_ui\_panels.py} and
\code{tests/test\_charts.py}.

Acceptance criterion: panels render without crashing; user inputs are
validated; expected outputs appear; and download and data-loading paths
behave consistently.

% ─────────────────────────────────────────────────────────────────────────────
\section{Backtesting Tests}

Required backtesting tests:

\begin{itemize}
  \item \textbf{BT\_01} Exception flag: $\mathrm{loss} > \mathrm{VaR}$ correctly identified
  \item \textbf{BT\_02} No-exception case (all losses below VaR)
  \item \textbf{BT\_03} All-exception case (all losses above VaR)
  \item \textbf{BT\_04} Expected exception count $= T \cdot (1-\alpha)$
  \item \textbf{BT\_05} Kupiec $\mathrm{LR}_{\mathrm{uc}}$ is finite and non-negative
  \item \textbf{BT\_06} Kupiec p-value in $[0,1]$
  \item \textbf{BT\_07} 95\% confidence backtest
  \item \textbf{BT\_08} 97.5\% confidence backtest
  \item \textbf{BT\_09} 99\% confidence backtest
  \item \textbf{BT\_10} Basel traffic-light zone correctly classified
  \item \textbf{BT\_11} Exception severity table
  \item \textbf{BT\_12} Christoffersen independence LR finite and non-negative
\end{itemize}

Lecture~5 emphasizes that a backtest must examine not only the
\emph{frequency} of exceptions but also their \emph{clustering} and
\emph{behavior across confidence levels}.  The repository includes
Christoffersen-style diagnostics; this test plan explicitly requires
coverage of those diagnostics.

% ─────────────────────────────────────────────────────────────────────────────
\section{Data Validation Tests}

\begin{itemize}
  \item \textbf{DATA\_01} Price series has strictly positive prices
  \item \textbf{DATA\_02} No duplicate dates after cleaning
  \item \textbf{DATA\_03} Missing-data report is generated for all-NaN columns
  \item \textbf{DATA\_04} Stale-price run of $\geq 10$ days is detected
  \item \textbf{DATA\_05} Extreme return outliers are visible (no silent clipping)
  \item \textbf{DATA\_06} Aligned histories share a common date index
  \item \textbf{DATA\_07} Insufficient lookback raises with a descriptive message
  \item \textbf{DATA\_08} CSV with a bad date column raises
  \item \textbf{DATA\_09} CSV with a missing price column raises
  \item \textbf{DATA\_10} Data-proxy caveat is documented in the test plan and software design
\end{itemize}

Lecture~5 explicitly warns that poor data quality leads to poor model
outputs.  Data validation therefore belongs in the test plan rather
than only in the software-design document.

% ─────────────────────────────────────────────────────────────────────────────
\section{Coverage Plan}

The coverage target is the highest achievable with the no-network unit
suite.  Streamlit UI branch paths are excluded because they require a
live browser context.  All other \code{src/} modules are expected to
reach or exceed 95\% statement coverage.  The acceptance command is:

\begin{shellcode}
python -m pytest tests/ \
  --cov=src \
  --cov-report=term-missing \
  --cov-report=html:submission/coverage_report \
  --cov-report=xml:submission/coverage_report/coverage.xml \
  --ignore=tests/integration_test.py \
  --ignore=tests/integration_test_formula_sheet.py
\end{shellcode}

Coverage steps:
\begin{enumerate}
  \item Run coverage and identify all missing lines.
  \item Add tests for missing branches wherever feasible.
  \item Use \code{pragma: no cover} only for genuinely unreachable
    defensive code.
  \item Save terminal coverage output and \code{htmlcov/index.html}.
  \item Include coverage output and missing-line discussion in Deliverable~5.
\end{enumerate}

% ─────────────────────────────────────────────────────────────────────────────
\section{Acceptance Criteria}

The software satisfies the planned test standard if:

\begin{enumerate}
  \item All required no-network unit tests pass.
  \item Analytical golden tests match expected values within stated tolerances.
  \item Homework-derived fixtures pass within stated tolerances.
  \item Invalid inputs produce controlled failures rather than silent bad outputs.
  \item Backtesting diagnostics return finite, interpretable results.
  \item UI smoke tests pass.
  \item Integration workflows either pass or have explicitly documented, understood failures.
  \item Statement coverage across \code{src/} reaches $\geq 95\%$ for all
    non-UI modules; Streamlit panel branches requiring a live browser session
    are excluded from the hard target and documented as justified gaps.
  \item Benchmark comparisons (Section~\ref{sec:mrm-benchmark}) confirm
    that relative differences between methods are within documented bounds.
\end{enumerate}

% ─────────────────────────────────────────────────────────────────────────────
\section{Known Untested Areas}

A sound test plan acknowledges residual validation risk honestly:

\begin{itemize}
  \item Some extension modules have historically lower coverage than the
    core market-risk engine.
  \item External option-calculator benchmarking may remain partially manual
    unless automated.
  \item Live market-data integration is inherently less deterministic than
    no-network unit tests.
  \item Duplicate-date or stale-price handling is not as extensively tested
    as basic missing-data and positivity checks.
  \item Coverage target may not reach 100\% without additional branch-specific
    tests in CDS, hazard absolute-shock branches, and selected UI/regulatory
    paths.
\end{itemize}

% ─────────────────────────────────────────────────────────────────────────────
\section{Appendix: Behavioural and Convergence Tests}

\subsection{Behavioural Tests}

\begin{itemize}
  \item \textbf{BEH\_01} Call price increases with spot
  \item \textbf{BEH\_02} Put price decreases with spot
  \item \textbf{BEH\_03} Option price increases with volatility
  \item \textbf{BEH\_04} Call price decreases with strike
  \item \textbf{BEH\_05} Put price increases with strike
  \item \textbf{BEH\_06} VaR increases with notional
  \item \textbf{BEH\_07} VaR increases with volatility
  \item \textbf{BEH\_08} CDS spread increases with hazard rate
  \item \textbf{BEH\_09} CVA increases with expected exposure
  \item \textbf{BEH\_10} CVA decreases with recovery rate
  \item \textbf{BEH\_11} Capital ratio decreases after stress losses
  \item \textbf{BEH\_12} Stressed capital-path minimum is correctly selected
\end{itemize}

\subsection{Convergence and Stability Tests}

\begin{itemize}
  \item \textbf{CONV\_01} One-asset MC VaR converges toward exact lognormal VaR as $n_{\mathrm{sims}}$ increases
  \item \textbf{CONV\_02} One-asset MC ES converges toward exact lognormal ES as $n_{\mathrm{sims}}$ increases
  \item \textbf{CONV\_03} Finite-difference delta converges to analytical Black-Scholes delta
  \item \textbf{CONV\_04} Too-small bump demonstrates cancellation risk or is explicitly avoided
  \item \textbf{CONV\_05} Rolling-window VaR sensitivity across window lengths (2y vs.\ 5y)
  \item \textbf{CONV\_06} EWMA $\lambda$ sensitivity
\end{itemize}

\subsection{Traceability to Repository Test Files}

\begin{itemize}
  \item \code{tests/test\_backend.py}
  \item \code{tests/test\_backtest\_extensions.py}
  \item \code{tests/test\_course\_validation.py}
  \item \code{tests/test\_homework\_cases.py}
  \item \code{tests/test\_lognormal.py}
  \item \code{tests/test\_credit.py}
  \item \code{tests/test\_credit\_service.py}
  \item \code{tests/test\_cva\_mitigants.py}
  \item \code{tests/test\_counterparty\_mitigation.py}
  \item \code{tests/test\_market\_data.py}
  \item \code{tests/test\_config\_and\_validation.py}
  \item \code{tests/test\_ui\_panels.py}
  \item \code{tests/test\_charts.py}
  \item \code{tests/test\_regulatory.py}
  \item \code{tests/test\_dfast\_pathing.py}
  \item \code{tests/test\_balance\_sheet.py}
  \item \code{tests/test\_coverage\_gaps.py}
  \item \code{tests/test\_strict\_numerics.py}
  \item \code{tests/test\_es\_confidence\_split.py}
  \item \code{tests/integration\_test.py}
  \item \code{tests/integration\_test\_formula\_sheet.py}
\end{itemize}

\end{document}
